% TOP_PROBLEM: {"id": 4600043, "title": "Growth of jointly periodic points I", "category_id": 11, "category": {"id": 11, "name": "dynamical_systems", "display_name": "Dynamical Systems", "description": "Problems about long-term behavior of deterministic systems, Hamiltonian dynamics, and periodic orbits.", "slug": "dynamical-systems", "order_index": 11, "created_at": "2026-07-31T15:26:25.671Z"}, "status": "open", "problem_number": "AMR-045-0043", "set_id": 13}
% FIRST_POSTED: 2026-10-01
% ENTRY_KIND: statement_only
% UnsolvedMath ID 4600043; source catalogue code AMR-045-0043
% !TeX program = lualatex
% Definition and sources only; this file is not an LLM solution attempt.
\documentclass[11pt,a4paper]{article}
\usepackage[margin=25mm]{geometry}
\usepackage{fontspec}
\setmainfont{lmroman10-regular.otf}[BoldFont=lmroman10-bold.otf,
 ItalicFont=lmroman10-italic.otf,BoldItalicFont=lmroman10-bolditalic.otf]
\usepackage{amsmath,amssymb,mathtools}
\usepackage{unicode-math}
\setmathfont{latinmodern-math.otf}
\AtBeginDocument{\let\setminus\smallsetminus}
\usepackage{xurl,hyperref}
\hypersetup{colorlinks=true,urlcolor=blue}
\setcounter{secnumdepth}{0}
\setlength{\parindent}{0pt}
\setlength{\parskip}{5pt}
\setlength{\emergencystretch}{3em}
\begin{document}
\section*{Growth of jointly periodic points I}\label{4600043}
{\small UnsolvedMath ID 4600043 · AMR-045-0043 · Dynamical Systems · AMR Open Problem Lists}
\subsection{Definitions and mathematical statement}
For $N\ge2$ let $X_N=\{0,\ldots,N-1\}^{\mathbb Z}$ and
$(\sigma_Nx)_j=x_{j+1}$. Let $F:X_N\to X_N$ be a surjective
cellular automaton, that is, a continuous map commuting with
$\sigma_N$, or equivalently a map specified by a finite local
rule. For each $k\ge1$ put
\[
 P_k(F)=\#\{x\in X_N:\sigma_N^k x=x
                  \text{ and }F^m x=x\text{ for some }m\ge1\},
 \qquad \nu(F)=\limsup_{k\to\infty}P_k(F)^{1/k}.
\]
The temporal period $m$ may depend on $x$ and is not bounded
in this definition. Spatial periods dividing $k$ are counted,
not just configurations of least spatial period $k$.

Must every such onto automaton satisfy $\nu(F)>1$?
Equivalently, for each $F$ must there be some $c>1$ for which
$P_k(F)\ge c^k$ for infinitely many $k$? The constant may
depend on the alphabet and the local rule.

The finite set $\operatorname{Fix}(\sigma_N^k)$ contains
$N^k$ points and $F$ acts on it as a self-map. Here $P_k$
counts precisely the points on its directed cycles; eventual
cycle entrants are excluded. Consequently $1\le\nu(F)\le N$.
The problem seeks a strictly positive exponential rate, rather
than merely infinitely many jointly periodic configurations.
\subsection{Short English statement}
Do the cycle points of an onto cellular automaton on spatially periodic configurations always grow at a strictly exponential rate?
\subsection{Sources}
\begin{itemize}
\item[S1] ULAM.AI and the UnsolvedMath Contributors, supplied problems.json, record 4600043 (AMR-045-0043), AMR Open Problem Lists. Catalogue identity and source transcription; CC BY 4.0.
\url{https://huggingface.co/datasets/ulamai/UnsolvedMath}.
\item[S2] Boyle - Open Problems in Symbolic Dynamics (2008), Problem/Question/Conjecture 25.3. Original source indexed by AMR-045-0043.
\url{https://www.math.umd.edu/~mboyle/papers/openfinalsub3nov2007.pdf}.
\end{itemize}
\end{document}
